\section{Aplicación en Python}

\begin{frame}{Ejemplo visual: lineal vs. RBF sobre una frontera curva}
\small
\begin{columns}[T,onlytextwidth]
\begin{column}{0.66\textwidth}
\IfFileExists{fig_svm_kernel_moons.png}{%
  \centering\includegraphics[width=\linewidth]{fig_svm_kernel_moons.png}
}{%
  \begin{ccdebox}\centering\small Ejecuta \texttt{python codigo\_demo\_svm.py} para generar esta comparación.\end{ccdebox}}
\end{column}
\begin{column}{0.30\textwidth}
\begin{contrastbox}{ccdeblue}{Kernel lineal}
\scriptsize
Solo puede construir un hiperplano en el espacio original. Con \texttt{make\_moons} queda sesgado.
\end{contrastbox}
\vspace{0.11cm}
\begin{contrastbox}{ccdeaccent}{Kernel RBF}
\scriptsize
Construye una separación curvada mediante similitudes locales en el espacio implícito.
\end{contrastbox}
\vspace{0.11cm}
\begin{keybox}[fontupper=\scriptsize]
La flexibilidad adicional debe validarse fuera de muestra.
\end{keybox}
\end{column}
\end{columns}
\source{Figura generada por \texttt{codigo\_demo\_svm.py}, semilla 42.}
\end{frame}

% ================================================================
\begin{frame}[fragile]{Python 1/3: pipeline y separación train/test}
\scriptsize
\begin{columns}[T,onlytextwidth]
\begin{column}{0.64\textwidth}
\begin{lstlisting}[style=ccdepython]
from sklearn.datasets import make_moons
from sklearn.model_selection import train_test_split
from sklearn.pipeline import make_pipeline
from sklearn.preprocessing import StandardScaler
from sklearn.svm import SVC

X, y = make_moons(
    n_samples=600, noise=0.25, random_state=42
)
X_train, X_test, y_train, y_test = train_test_split(
    X, y, test_size=0.30,
    stratify=y, random_state=42
)

model = make_pipeline(
    StandardScaler(),
    SVC(kernel="rbf", C=10, gamma="scale")
)
model.fit(X_train, y_train)
\end{lstlisting}
\end{column}
\begin{column}{0.32\textwidth}
\begin{keybox}
\textbf{Puntos clave}
\begin{itemize}
  \item Escalamiento dentro del pipeline.
  \item \texttt{stratify} conserva proporciones.
  \item Semilla fija para reproducibilidad.
  \item RBF para una frontera no lineal.
\end{itemize}
\end{keybox}
\end{column}
\end{columns}
\end{frame}

% ================================================================
\begin{frame}[fragile]{Python 2/3: validación cruzada de $C$, $\gamma$ y kernel}
\scriptsize
\begin{columns}[T,onlytextwidth]
\begin{column}{0.67\textwidth}
\begin{lstlisting}[style=ccdepython]
from sklearn.model_selection import GridSearchCV

pipe = make_pipeline(StandardScaler(), SVC())
param_grid = [
    {
        "svc__kernel": ["linear"],
        "svc__C": [0.1, 1, 10, 100],
    },
    {
        "svc__kernel": ["rbf"],
        "svc__C": [0.1, 1, 10, 100],
        "svc__gamma": ["scale", 0.1, 1, 10],
    },
]

search = GridSearchCV(
    pipe, param_grid, scoring="roc_auc",
    cv=5, n_jobs=-1, refit=True
)
search.fit(X_train, y_train)
\end{lstlisting}
\end{column}
\begin{column}{0.29\textwidth}
\begin{keybox}
\textbf{Por qué así}
\begin{itemize}
  \item La grilla vive solo en training.
  \item Escalamiento se re-estima en cada fold.
  \item Se compara lineal vs. RBF.
  \item ROC-AUC guía la selección.
\end{itemize}
\end{keybox}
\end{column}
\end{columns}
\end{frame}

% ================================================================
\begin{frame}[fragile]{Python 3/3: evaluación fuera de muestra}
\scriptsize
\begin{columns}[T,onlytextwidth]
\begin{column}{0.65\textwidth}
\begin{lstlisting}[style=ccdepython]
from sklearn.metrics import (
    accuracy_score, confusion_matrix, roc_auc_score
)

best = search.best_estimator_
pred = best.predict(X_test)
score = best.decision_function(X_test)

print(search.best_params_)
print("accuracy =", accuracy_score(y_test, pred))
print("roc_auc =", roc_auc_score(y_test, score))
print(confusion_matrix(y_test, pred))
print(
    "support vectors =",
    best.named_steps["svc"].n_support_
)
\end{lstlisting}
\end{column}
\begin{column}{0.31\textwidth}
\begin{contrastbox}{ccdeaccent}{Resultado reproducible}
\scriptsize
Mejor grilla: RBF, $C=10$, $\gamma=\texttt{scale}$.\\[0.10cm]
ROC-AUC CV: \textbf{0.9724}.\\
Accuracy test: \textbf{0.9167}.\\
ROC-AUC test: \textbf{0.9817}.\\
Support vectors: $42+39=81$.
\end{contrastbox}
\end{column}
\end{columns}
\source{Ejecución local de \texttt{codigo\_demo\_svm.py}; scikit-learn, semilla 42.}
\end{frame}

% ================================================================
\begin{frame}{Lineal vs. RBF: resultado del mismo experimento}
\begin{columns}[T,onlytextwidth]
\begin{column}{0.58\textwidth}
\vspace{1.55cm}
\renewcommand{\arraystretch}{1.45}
\small
\begin{tabular}{lccc}
\toprule
\rowcolor{ccdenavy}
\textcolor{white}{\textbf{Modelo}} &
\textcolor{white}{\textbf{Accuracy}} &
\textcolor{white}{\textbf{ROC-AUC}} &
\textcolor{white}{\textbf{SV}}\\
\midrule
SVM lineal & 0.8556 & 0.9383 & 141\\
\rowcolor{ccdeice!50}
SVM RBF & \textbf{0.9333} & \textbf{0.9775} & 75\\
\bottomrule
\end{tabular}
\end{column}
\begin{column}{0.38\textwidth}
\begin{keybox}
\textbf{Lectura}
\begin{itemize}
  \item La estructura de \texttt{make\_moons} es no lineal.
  \item El RBF reduce el sesgo respecto al hiperplano lineal.
  \item Menos support vectors no implica siempre mejor modelo, pero aquí la frontera RBF es además más parsimoniosa en puntos activos.
\end{itemize}
\end{keybox}
\end{column}
\end{columns}
\source{Ejecución de \texttt{codigo\_demo\_svm.py}, semilla 42.}
\end{frame}

% ================================================================
